\documentclass[12pt,a4paper]{article}

\usepackage[utf8]{inputenc}
\usepackage[T1]{fontenc}
\usepackage[french]{babel}
\usepackage{geometry}
\usepackage{graphicx}
\usepackage{listings}

\geometry{margin=2.5cm}
\lstset{
  language=Python,
  basicstyle=\ttfamily\small,
  frame=single,
  breaklines=true,
  showstringspaces=false,
  literate={é}{{\'e}}1 {à}{{\`a}}1 {’}{{'}}1
}

\newcommand{\rapportcapture}[2]{%
  \begin{figure}[htbp]
    \centering
    \IfFileExists{#1}{%
      \includegraphics[width=0.9\linewidth]{#1}%
    }{%
            \fbox{\parbox[c][4cm][c]{0.82\linewidth}{\centering #2}}
    }
  \end{figure}
}

\title{TP2}
\author{}
\date{}

\begin{document}

\maketitle

\section{Exercice 1 : Tests unitaires}

\paragraph{tests/unit/test\_livraison.py}
\begin{lstlisting}
import pytest

from app.livraison import frais_livraison


def test_livraison_panier_ordinaire():
    assert frais_livraison(120) == 25


def test_livraison_panier_299():
    assert frais_livraison(299) == 25


def test_livraison_panier_300():
    assert frais_livraison(300) == 0


def test_livraison_panier_301():
    assert frais_livraison(301) == 0


def test_livraison_panier_zero():
    assert frais_livraison(0) == 25


def test_livraison_montant_negatif():
    with pytest.raises(ValueError):
        frais_livraison(-1)
\end{lstlisting}

\rapportcapture
  {screenshots/fig1_output_exercice1.png}
  {Fig. 1 : Output exercice 1}

\section{Exercice 2 : Nettoyer un texte avant son analyse}

\paragraph{app/texte.py}
\begin{lstlisting}
def nettoyer_texte(texte):
  if not isinstance(texte, str):
    raise TypeError("Texte attendu")
  return " ".join(texte.lower().split())
\end{lstlisting}

\paragraph{tests/test\_texte.py}
\begin{lstlisting}
import pytest

from app.texte import nettoyer_texte


@pytest.mark.parametrize(
  "entree, attendu",
  [
    ("  BONJOUR   monde  ", "bonjour monde"),
    (" déjà propre ", "déjà propre"),
    ("", ""),
    ("   ", ""),
    (" PYTHON ", "python"),
  ],
)
def test_nettoyer_texte(entree, attendu):
  assert nettoyer_texte(entree) == attendu


@pytest.mark.parametrize("entree", [123, None])
def test_nettoyer_texte_entree_non_textuelle(entree):
  with pytest.raises(TypeError):
    nettoyer_texte(entree)


def test_nettoyer_texte_deja_nettoye_ne_change_pas():
  texte_nettoye = nettoyer_texte("  BONJOUR   monde  ")
  assert nettoyer_texte(texte_nettoye) == texte_nettoye
\end{lstlisting}

\rapportcapture
  {screenshots/fig2_output_exercice2.png}
  {Fig. 2 : Output exercice 2}

\section{Exercice 3 : Fixtures et tests d'intégration}

\paragraph{app/reservation.py}
\begin{lstlisting}
class Stock:
  def __init__(self, quantite):
    self.quantite = quantite

  def retirer(self, quantite):
    if quantite <= 0:
      raise ValueError("Quantité positive requise")
    if quantite > self.quantite:
      return False
    self.quantite -= quantite
    return True


class Reservation:
  def __init__(self, stock):
    self.stock = stock

  def reserver(self, quantite):
    if self.stock.retirer(quantite):
      return "confirmee"
    return "refusee"
\end{lstlisting}

\paragraph{tests/test\_reservation.py}
\begin{lstlisting}
import pytest

from app.reservation import Reservation, Stock


@pytest.fixture
def stock():
  return Stock(5)


def test_reservation_deux_unites_confirmee_et_stock_reduit(stock):
  reservation = Reservation(stock)
  assert reservation.reserver(2) == "confirmee"
  assert stock.quantite == 3


def test_reservation_six_unites_refusee_et_stock_inchange(stock):
  reservation = Reservation(stock)
  assert reservation.reserver(6) == "refusee"
  assert stock.quantite == 5


def test_reservation_cinq_unites_confirmee_et_stock_vide(stock):
  reservation = Reservation(stock)
  assert reservation.reserver(5) == "confirmee"
  assert stock.quantite == 0


def test_reservation_zero_unite_leve_value_error_et_stock_inchange(stock):
  reservation = Reservation(stock)
  with pytest.raises(ValueError):
    reservation.reserver(0)
  assert stock.quantite == 5
\end{lstlisting}

\rapportcapture
  {screenshots/fig3_output_exercice3.png}
  {Fig. 3 : Output exercice 3}

\section{Exercice 4 : Isolation des dépendances}

\paragraph{app/moderation.py}
\begin{lstlisting}
def decider(modele, message):
  classe = modele.predict(message)
  if classe == "spam":
    return "bloquer"
  return "autoriser"
\end{lstlisting}

\paragraph{tests/test\_moderation.py}
\begin{lstlisting}
from unittest.mock import Mock

import pytest

from app.moderation import decider


def test_decider_bloque_un_spam():
  modele = Mock()
  modele.predict.return_value = "spam"
  assert decider(modele, "Offre miracle") == "bloquer"
  modele.predict.assert_called_once_with("Offre miracle")


def test_decider_autorise_un_message_normal():
  modele = Mock()
  modele.predict.return_value = "normal"
  assert decider(modele, "Bonjour") == "autoriser"
  modele.predict.assert_called_once_with("Bonjour")


def test_decider_propage_erreur_du_modele():
  modele = Mock()
  modele.predict.side_effect = RuntimeError("Modele indisponible")
  with pytest.raises(RuntimeError, match="Modele indisponible"):
    decider(modele, "Offre miracle")
\end{lstlisting}

\rapportcapture
  {screenshots/fig4_output_exercice4.png}
  {Fig. 4 : Output exercice 4}

\end{document}